\chapter{Complex Numbers and Phasors}
\label{ch:complex}
Much of the analysis in this book---and almost all of the steady-state circuit
work in Parts~III and~IV---rests on two ideas from complex algebra: a complex number, and
a \emph{phasor}, the complex number that stands in for a steady sinusoid. This
short chapter is a self-contained refresher on both. It stays generic: a phasor
is a tool for \emph{any} sinusoidally driven linear system, mechanical or
electrical, and nothing here assumes a circuit. The circuit-specific
connections---impedance, admittance, and complex power---wait until \cref{ch:ac}.

\section{Complex Number Refresher}
\label{sec:complexrefresher}
A complex number in \textbf{rectangular form} is \(z=a+\jj b\), with real part
\(a\) and imaginary part \(b\). Addition and subtraction are component-wise:
\[
(a+\jj b)+(c+\jj d)=(a+c)+\jj(b+d).
\]
In \textbf{polar form} \(z=r\angle\theta\) with
\[
r=\sqrt{a^2+b^2},\qquad \theta=\operatorname{atan2}(b,a),
\]
and multiplication and division are easiest there:
\[
(r_1\angle\theta_1)(r_2\angle\theta_2)=r_1r_2\angle(\theta_1+\theta_2),
\qquad
\frac{r_1\angle\theta_1}{r_2\angle\theta_2}=\frac{r_1}{r_2}\angle(\theta_1-\theta_2).
\]
The two forms are joined by \textbf{Euler's formula},
\(\ee^{\jj\theta}=\cos\theta+\jj\sin\theta\), the bridge between sinusoids and
complex algebra.

\begin{mistakebox}
\(\tan^{-1}(b/a)\) alone can land the angle in the wrong quadrant. Use
\(\operatorname{atan2}(b,a)\), or reason from the signs of both coordinates.
\end{mistakebox}

\begin{workedbox}[: converting between forms]
Convert \(z=3+\jj4\) to polar: \(r=\sqrt{9+16}=5\),
\(\theta=\operatorname{atan2}(4,3)\approx53.1^\circ\), so \(z=5\angle53.1^\circ\).
Going back, \(10\angle{-30^\circ}=10\cos(-30^\circ)+\jj\,10\sin(-30^\circ)
\approx 8.66-\jj5\).
\end{workedbox}

\begin{exambox}
\textbf{Remember this.} \emph{Use rectangular form to add, polar form to
multiply.} Nearly every impedance calculation in this book is some sequence of
those two moves, and choosing the wrong form is what turns a one-line problem into
a page of algebra.
\end{exambox}

\section{Why Complex Numbers Appear: Phasors}
Drive a linear system with a steady sinusoid and, once transients die out, every
voltage, current, force, or flow in it is a sinusoid at the \emph{same}
frequency---differing only in amplitude and phase. A \textbf{phasor} captures
exactly those two numbers, amplitude and phase, as a single complex number, so
that the shared \(\cos\omega t\) can be captured instead of written out everywhere.

Concretely, a sinusoid is written as the real part of a rotating complex
exponential:
\[
x(t)=\sqrt{2}\,\Re\{X\,\ee^{\jj\omega t}\}.
\]
Here \(X\) is the \textbf{phasor}: a \emph{complex} constant whose magnitude is
the signal's \textbf{root-mean-square (RMS)} amplitude and whose angle is its
phase, \(X=|X|\angle\phi\). It does not depend on time; all the time dependence
sits in the universal factor \(\ee^{\jj\omega t}\), which simply spins the constant
\(X\) around the origin---and the sinusoid is that spinning point's shadow on the
real axis (\cref{fig:phasor}). The \(\sqrt2\) is a deliberate choice of scale, and
the next section explains it.

Two angle symbols are used consistently from here on, and both are listed in the
Notation table: \(\phi\) is a \emph{signal's phase}---the angle of a phasor---while
\(\theta\) is the generic polar-form angle \(z=r\angle\theta\) of
\cref{sec:complexrefresher}, which is also what the impedance angle of
\cref{ch:ac} is called.

\section{Root Mean Square: the Amplitude That Measures Power}
\label{sec:rms}
This book's phasor magnitudes are \textbf{RMS} values by default, and that
convention is used from the study guide onward---so it is worth deriving rather
than announcing, because the \(\sqrt2\) above is its direct consequence.

The problem RMS solves is that a sinusoid has no single obvious ``size.'' Its
average over a cycle is zero, and its peak is reached only instantaneously. What
usually matters physically is neither: it is how much \emph{power} the signal
delivers. So define the amplitude that answers that question, and let everything
else follow.

Take a sinusoid \(x(t)=X_p\cos\omega t\) driving a resistance \(R\). Instantaneous
power is \(p=x^2/R\), so the average over one period \(T=2\pi/\omega\) is
\[
\bar P=\frac{1}{T}\int_0^T\frac{X_p^2\cos^2\omega t}{R}\,\dd t
=\frac{X_p^2}{R}\cdot\frac{1}{T}\int_0^T\cos^2\omega t\,\dd t .
\]
The remaining integral is the whole calculation, and the half-angle identity
\(\cos^2\theta=\tfrac12(1+\cos2\theta)\) does it in one line: the \(\cos2\omega t\)
term integrates to zero over a whole number of cycles, leaving just the \(\tfrac12\).
Hence
\[
\bar P=\frac{X_p^2}{2R}
=\frac{\bigl(X_p/\sqrt2\bigr)^2}{R} .
\]
Read the right-hand form: the average power is what a \emph{constant} value
\(X_p/\sqrt2\) would dissipate in the same resistor. That constant is the
\textbf{root-mean-square} value---literally the square \emph{root} of the
\emph{mean} of the \emph{square},
\[
X_{\RMS}=\sqrt{\frac{1}{T}\int_0^T x^2(t)\,\dd t}
\qquad\Longrightarrow\qquad
\boxed{\;X_{\RMS}=\frac{X_p}{\sqrt2}\approx0.707\,X_p\quad\text{(sinusoid only)}\;}
\]
Two consequences. First, using RMS magnitudes makes the DC power formulas work
unchanged on AC---\(P=V_{\RMS}^2/R=I_{\RMS}^2R\), with no stray factors of two,
which is exactly why the convention is standard in power and radio work. Second,
since the peak is \(\sqrt2\) times the RMS value, recovering the time-domain
waveform from an RMS phasor requires putting that \(\sqrt2\) back---which is the
\(\sqrt2\) in \(x(t)=\sqrt2\,\Re\{X\ee^{\jj\omega t}\}\).

\begin{mistakebox}
\(X_{\RMS}=X_p/\sqrt2\) holds \emph{for a sinusoid}, and for nothing else. The
definition---root of the mean of the square---is general, but the \(0.707\) factor
is not: a square wave of peak \(X_p\) has \(X_{\RMS}=X_p\), and a triangle wave has
\(X_p/\sqrt3\). \Cref{ch:ac} needs this distinction when a waveform is not
sinusoidal.
\end{mistakebox}

\begin{workedbox}[: household mains]
A \SI{120}{\volt} outlet is quoted in RMS, because that is the number that predicts
heating. Its peak and peak-to-peak values are therefore
\[
V_p=\sqrt2\,(120)\approx\SI{169.7}{\volt},
\qquad
V_{pp}=2V_p\approx\SI{339.4}{\volt} .
\]
Insulation has to withstand the peak; a heater dissipates according to the RMS
value. Both numbers describe one waveform, and quoting the wrong one is off by
\(41\%\) or \(183\%\).
\end{workedbox}

\begin{figure}[htbp]
\centering
\includegraphics[width=0.86\textwidth]{phasor_unit_circle}
\caption{The phasor picture. \emph{Left:} the complex constant
\(X=|X|\angle\phi\) shown in both polar (magnitude and angle) and rectangular
(\(a+\jj b\)) form. \emph{Right:} multiplying by \(\ee^{\jj\omega t}\) spins it
counter-clockwise at rate \(\omega\); the real-axis projection traces the
time-domain sinusoid \(x(t)=\sqrt{2}\,\Re\{X\ee^{\jj\omega t}\}\).}
\label{fig:phasor}
\end{figure}

\subsection{Where the Formula Comes From}
The expression \(x(t)=\sqrt2\,\Re\{X\ee^{\jj\omega t}\}\) is just Euler's formula
in disguise. Write the phasor in polar form, \(X=|X|\ee^{\jj\phi}\). Then
\[
X\ee^{\jj\omega t}=|X|\ee^{\jj(\omega t+\phi)}
=|X|\bigl[\cos(\omega t+\phi)+\jj\sin(\omega t+\phi)\bigr],
\]
using Euler's formula on the last step. Taking the real part keeps the cosine,
\[
\Re\{X\ee^{\jj\omega t}\}=|X|\cos(\omega t+\phi),
\qquad\text{so}\qquad
x(t)=\sqrt2\,|X|\cos(\omega t+\phi).
\]
So the phasor's magnitude \(|X|\) and angle \(\phi\) are precisely the amplitude
and phase of the sinusoid. The \(\sqrt2\) is a bookkeeping choice: this book uses
\textbf{RMS} phasor magnitudes by default, and a sinusoid of RMS value \(|X|\) has
peak value \(\sqrt2\,|X|\), which is where the factor comes from. (Some texts omit
the \(\sqrt2\) and let \(|X|\) be the peak amplitude; watch the convention.)

\begin{controlsbox}
A phasor is the \(s=\jj\omega\) snapshot of the more general story in
\cref{ch:linsys,ch:splane}. There, a linear system's response is tracked at any
complex frequency \(s=\sigma+\jj\omega\); the phasor is the special case of a
\emph{pure} sinusoid, \(\sigma=0\), evaluated on the imaginary axis.
\end{controlsbox}

The single most useful consequence: differentiation in time becomes
multiplication by \(\jj\omega\),
\[
\frac{\dd}{\dd t}\bigl(X\ee^{\jj\omega t}\bigr)=\jj\omega X\ee^{\jj\omega t}.
\]
Every derivative in a linear differential equation turns into a factor of
\(\jj\omega\), so the differential equations of a sinusoidally driven system
collapse into ordinary algebra in the phasors. That is the whole reason phasors
earn their keep, and it is the machinery \cref{ch:ac} uses to define impedance for
real circuits.
